% =============================================================================
% embedded-grundlagen.tex – Systematische Einführung
% =============================================================================

\section{Embedded-Grundlagen – Systematische Einführung}

\subsection{Begriffsrahmen und Zielsetzung}

Embedded-Systeme unterscheiden sich fundamental von Desktop-Software. Die \glqq API\grqq{} ist physisch – Pins, Spannungspegel, Timing. Fehler führen nicht zu Exceptions, sondern zu falschen Signalen, zerstörter Hardware oder unerklärlichem Verhalten.

\begin{table}[htbp]
\centering
\begin{tabular}{@{}ll@{}}
\toprule
\textbf{Begriff} & \textbf{Definition} \\
\midrule
GPIO            & General Purpose Input/Output – frei konfigurierbarer Pin \\
HAL             & Hardware Abstraction Layer – Schicht zwischen Software und Registern \\
Active-Low      & Signal ist aktiv bei Pegel 0 (GND) \\
Active-High     & Signal ist aktiv bei Pegel 1 (VCC) \\
Debouncing      & Entprellung – Filterung mechanischer Prellvorgänge \\
SPI             & Serial Peripheral Interface – synchroner Datenbus \\
Schieberegister & IC, das parallele und serielle Daten konvertiert \\
\bottomrule
\end{tabular}
\caption{Begriffsrahmen Embedded-Grundlagen}
\label{tab:grundlagen-begriffe}
\end{table}

\textbf{Ziel:} Du verstehst die physikalischen Grundlagen, die hinter dem Selection-Panel-Code stehen.

% -----------------------------------------------------------------------------
\subsection{Hardware ist die Schnittstelle}

\textbf{Problem:} In Desktop-Software rufst du Funktionen auf und bekommst Rückgabewerte. In Embedded-Systemen kommunizierst du mit der Außenwelt über elektrische Signale.

\begin{figure}[htbp]
\centering
\begin{tikzpicture}[
    layer/.style={draw, fill=CodeBg, minimum width=11cm, minimum height=0.8cm,
                  font=\small, rounded corners=2pt},
    arr/.style={-{Stealth[length=2mm]}, thick, color=ArduinoDark}
]
\node[layer] (app) at (0,4) {Anwendungslogik (Selection, Debounce)};
\node[layer] (drv) at (0,3) {Treiber (CD4021, HC595)};
\node[layer] (hal) at (0,2) {HAL (SPI-Transfer, GPIO-Zugriff)};
\node[layer] (reg) at (0,1) {Register (ESP32-S3 Peripherie)};
\node[layer] (pin) at (0,0) {Physische Pins (\SI{3,3}{\volt} / GND)};
\draw[arr] (app.south) -- (drv.north);
\draw[arr] (drv.south) -- (hal.north);
\draw[arr] (hal.south) -- (reg.north);
\draw[arr] (reg.south) -- (pin.north);
\end{tikzpicture}
\caption{Schichten vom Code zur Hardware}
\label{fig:hardware-schichten}
\end{figure}

\begin{warnbox}[Konsequenz]
Ein Fehler auf jeder Ebene führt zu Fehlfunktion. Wenn der SPI-Modus falsch ist, liefert der CD4021 Müll – auch wenn dein Code syntaktisch korrekt ist.
\end{warnbox}

% -----------------------------------------------------------------------------
\subsection{Spannungspegel – Die Sprache der Hardware}

\textbf{Grundprinzip:} Digitale Signale kennen zwei Zustände: HIGH und LOW.

\begin{table}[htbp]
\centering
\begin{tabular}{@{}lll@{}}
\toprule
\textbf{Zustand} & \textbf{Spannung (\SI{3,3}{\volt} Logik)} & \textbf{Bedeutung} \\
\midrule
LOW         & \SI{0}{\volt} – \SI{0,8}{\volt}   & Logisch 0 \\
HIGH        & \SI{2,0}{\volt} – \SI{3,3}{\volt} & Logisch 1 \\
Undefiniert & \SI{0,8}{\volt} – \SI{2,0}{\volt} & Verboten! \\
\bottomrule
\end{tabular}
\caption{Spannungspegel bei \SI{3,3}{\volt}-Logik}
\label{tab:spannungspegel}
\end{table}

\textbf{Warum Pullups/Pulldowns?}

Ein offener Pin (nicht verbunden) hat keinen definierten Pegel – er \glqq schwebt\grqq{} und nimmt zufällige Werte an. Ein Pullup-Widerstand zieht den Pin auf HIGH, solange nichts anderes ihn auf LOW zieht.

\begin{figure}[htbp]
\centering
\begin{tikzpicture}[
    comp/.style={font=\small},
    wire/.style={thick}
]
% VCC
\node[comp] (vcc) at (0,3) {VCC (\SI{3,3}{\volt})};
\draw[wire] (0,2.7) -- (0,2.3);

% Resistor (simple rectangle)
\draw[wire] (-0.15,2.3) rectangle (0.15,1.5);
\node[comp, right=0.3cm] at (0.15,1.9) {\SI{10}{\kilo\ohm} (Pullup)};

% Junction and wire to pin
\draw[wire] (0,1.5) -- (0,1);
\filldraw (0,1) circle (2pt);
\draw[wire] (0,1) -- (2,1);
\node[comp, right] at (2,1) {zum Pin};

% Switch (simple representation)
\draw[wire] (0,1) -- (0,0.3);
\draw[wire] (-0.3,0.3) -- (0.3,0) node[comp, right=0.2cm] {Taster};
\draw[wire] (0,-0.3) -- (0,-0.7);
\node[comp] at (0,-1) {GND};

% Labels
\node[comp, color=MutedText, align=left] at (4,0.5) {offen = HIGH\\gedrückt = LOW};
\end{tikzpicture}
\caption{Pull-up-Beschaltung für Taster}
\label{fig:pullup}
\end{figure}

\begin{infobox}[Anwendung im Projekt]
Die CD4021-Eingänge haben externe Pullups. Im Ruhezustand liest du HIGH (1), bei gedrücktem Taster LOW (0) – das ist Active-Low.
\end{infobox}

% -----------------------------------------------------------------------------
\subsection{Active-Low vs. Active-High}

\textbf{Problem:} Warum ist ein gedrückter Taster \glqq 0\grqq{} und eine leuchtende LED \glqq 1\grqq? Das scheint unintuitiv.

\begin{table}[htbp]
\centering
\begin{tabular}{@{}lllll@{}}
\toprule
\textbf{Typ} & \textbf{Aktiv bei} & \textbf{Inaktiv bei} & \textbf{Typischer Einsatz} \\
\midrule
Active-Low  & 0 (GND) & 1 (VCC) & Taster mit Pullup, Reset-Signale \\
Active-High & 1 (VCC) & 0 (GND) & LEDs, Enable-Signale \\
\bottomrule
\end{tabular}
\caption{Active-Low vs. Active-High}
\label{tab:active-low-high}
\end{table}

\textbf{Warum Active-Low für Taster?}

Mit Pullup ist der Ruhezustand definiert (HIGH). Der Taster muss nur nach GND schalten – das ist schaltungstechnisch einfacher und robuster gegen Störungen.

\textbf{Warum Active-High für LEDs?}

LEDs leuchten, wenn Strom fließt. Bei Active-High: Pin HIGH → Strom fließt durch LED nach GND → LED leuchtet.

\begin{center}
\code{Pin (HIGH) -{}- 330$\Omega$ -{}- LED -{}- GND}\\[0.5em]
\footnotesize\color{MutedText} \SI{3,3}{\volt} \hspace{1.5cm} $\sim$\SI{2}{\volt} \hspace{1cm} $\sim$\SI{0}{\volt}
\end{center}

\textbf{Konsequenz im Code:}

\begin{lstlisting}[style=arduino]
// Button lesen (Active-Low)
bool isPressed = (raw & mask) == 0;  // Gedrueckt wenn Bit = 0

// LED setzen (Active-High)
if (ledOn) {
    ledState |= mask;   // Bit = 1 -> LED an
} else {
    ledState &= ~mask;  // Bit = 0 -> LED aus
}
\end{lstlisting}

% -----------------------------------------------------------------------------
\subsection{Schieberegister – Pins multiplizieren}

\textbf{Problem:} Der ESP32-S3 hat begrenzte GPIO-Pins. Für 100 Buttons und 100 LEDs bräuchtest du 200 Pins – die gibt es nicht.

\textbf{Lösung:} Schieberegister wandeln zwischen parallelen und seriellen Daten um.

\paragraph{CD4021 – Parallel zu Seriell (Buttons lesen)}

\begin{figure}[htbp]
\centering
\begin{tikzpicture}[
    ic/.style={draw, fill=CodeBg, minimum width=5cm, minimum height=1.5cm,
               font=\small, rounded corners=2pt},
    arr/.style={-{Stealth[length=2mm]}, thick, color=ArduinoDark},
    pin/.style={font=\footnotesize}
]
% Input pins
\foreach \i in {0,...,7} {
    \node[pin] at ({-2.5+\i*0.7},2.5) {D\pgfmathparse{int(7-\i)}\pgfmathresult};
    \draw[arr] ({-2.5+\i*0.7},2.2) -- ({-2.5+\i*0.7},1.5);
}
\node[pin, color=MutedText] at (0,3) {← 8 Buttons parallel};

% IC
\node[ic] (cd4021) at (0,0.5) {CD4021 Schieberegister};

% Output
\draw[arr] (cd4021.east) -- ++(1.5,0) node[right, font=\small] {Serial Out → ESP32};
\node[pin, color=MutedText] at (4.5,-0.2) {(1 Pin)};
\end{tikzpicture}
\caption{CD4021: Parallel-zu-Seriell-Wandlung}
\label{fig:cd4021}
\end{figure}

\textbf{Ablauf:}
\begin{enumerate}
\item \textbf{PS/LOAD = HIGH:} Parallele Eingänge werden ins Register übernommen
\item \textbf{PS/LOAD = LOW:} Register ist bereit zum Auslesen
\item \textbf{8× CLK-Pulse:} Bits werden nacheinander auf Serial Out geschoben
\end{enumerate}

\textbf{Timing:}

\begin{table}[htbp]
\centering
\begin{tabular}{@{}llll@{}}
\toprule
\textbf{Phase} & \textbf{PS/LOAD} & \textbf{CLK} & \textbf{Dauer} \\
\midrule
Load  & HIGH & –      & \SI{1}{\micro\second} \\
Shift & LOW  & $\uparrow\downarrow$ × 8 & $\sim$\SI{16}{\micro\second} \\
\bottomrule
\end{tabular}
\caption{CD4021-Timing}
\label{tab:cd4021-timing}
\end{table}

\paragraph{74HC595 – Seriell zu Parallel (LEDs schreiben)}

\begin{figure}[htbp]
\centering
\begin{tikzpicture}[
    ic/.style={draw, fill=CodeBg, minimum width=5cm, minimum height=1.5cm,
               font=\small, rounded corners=2pt},
    arr/.style={-{Stealth[length=2mm]}, thick, color=ArduinoDark},
    pin/.style={font=\footnotesize}
]
% Input
\node[font=\small] at (-4,0.5) {ESP32 →};
\draw[arr] (-3,0.5) -- (-2.5,0.5);
\node[pin, color=MutedText] at (-3.5,0) {(1 Pin)};

% IC
\node[ic] (hc595) at (0,0.5) {74HC595 Schieberegister};

% Output pins
\foreach \i in {0,...,7} {
    \draw[arr] ({-2.5+\i*0.7},-0.5) -- ({-2.5+\i*0.7},-1.2);
    \node[pin] at ({-2.5+\i*0.7},-1.5) {Q\pgfmathparse{int(7-\i)}\pgfmathresult};
}
\node[pin, color=MutedText] at (0,-2) {← 8 LEDs parallel};
\end{tikzpicture}
\caption{74HC595: Seriell-zu-Parallel-Wandlung}
\label{fig:hc595}
\end{figure}

\textbf{Ablauf:}
\begin{enumerate}
\item \textbf{8× SRCLK-Pulse:} Bits werden ins Schieberegister geschoben
\item \textbf{LATCH-Puls (RCK):} Daten werden an die Ausgänge übernommen
\end{enumerate}

\begin{infobox}[Warum der LATCH?]
Ohne Latch würden die LEDs während des Schiebens flackern. Der Latch hält die alten Werte, bis alle neuen Bits da sind.
\end{infobox}

\paragraph{Kaskadierung – Mehr als 8 Kanäle}

Mehrere Schieberegister werden in Reihe geschaltet:

\begin{lstlisting}[style=wrap,numbers=none]
CD4021 (IC1) --> CD4021 (IC0) --> ESP32
   Buttons 9-16     Buttons 1-8      MISO

ESP32 --> HC595 (IC0) --> HC595 (IC1)
   MOSI     LEDs 1-8        LEDs 9-16
\end{lstlisting}

\begin{tipbox}[Konsequenz]
Mit 13 CD4021 und 13 HC595 steuerst du 100 Buttons und 100 LEDs über nur 5 Pins (SCK, MISO, MOSI, PS, LATCH).
\end{tipbox}

% -----------------------------------------------------------------------------
\subsection{SPI – Der Datenbus}

\textbf{Problem:} Wie überträgst du Daten zwischen ESP32 und Schieberegistern synchron und zuverlässig?

\textbf{Lösung:} SPI (Serial Peripheral Interface) – ein synchroner Bus mit Takt.

\paragraph{Signale}

\begin{table}[htbp]
\centering
\begin{tabular}{@{}lll@{}}
\toprule
\textbf{Signal} & \textbf{Richtung} & \textbf{Funktion} \\
\midrule
SCK   & Master → Slave & Taktsignal (synchronisiert Übertragung) \\
MOSI  & Master → Slave & Master Out, Slave In (Daten zum Slave) \\
MISO  & Slave → Master & Master In, Slave Out (Daten vom Slave) \\
CS/SS & Master → Slave & Chip Select (welcher Slave ist aktiv) \\
\bottomrule
\end{tabular}
\caption{SPI-Signale}
\label{tab:spi-signale}
\end{table}

\paragraph{SPI-Modi}

\textbf{Problem:} Wann genau wird ein Bit gelesen – bei steigender oder fallender Flanke? Das hängt vom SPI-Modus ab.

\begin{table}[htbp]
\centering
\begin{tabular}{@{}cccll@{}}
\toprule
\textbf{Mode} & \textbf{CPOL} & \textbf{CPHA} & \textbf{Clock Idle} & \textbf{Sample bei} \\
\midrule
0 & 0 & 0 & LOW  & Rising Edge \\
1 & 0 & 1 & LOW  & Falling Edge \\
2 & 1 & 0 & HIGH & Falling Edge \\
3 & 1 & 1 & HIGH & Rising Edge \\
\bottomrule
\end{tabular}
\caption{SPI-Modi}
\label{tab:spi-modi}
\end{table}

\textbf{Im Projekt:}

\begin{table}[htbp]
\centering
\begin{tabular}{@{}lcll@{}}
\toprule
\textbf{IC} & \textbf{Mode} & \textbf{Frequenz} & \textbf{Begründung} \\
\midrule
CD4021  & 1 & \SI{500}{\kilo\hertz} & Konservativ für Breadboard \\
74HC595 & 0 & \SI{1}{\mega\hertz}   & Standard \\
\bottomrule
\end{tabular}
\caption{SPI-Konfiguration im Projekt}
\label{tab:spi-projekt}
\end{table}

\begin{warnbox}[Konsequenz]
Wenn du den falschen Modus verwendest, liest du Müll oder die Daten sind um ein Bit verschoben.
\end{warnbox}

% -----------------------------------------------------------------------------
\subsection{Timing – Determinismus ist Pflicht}

\textbf{Problem:} Embedded-Software muss in definierten Zeitfenstern reagieren. Wenn du eine Deadline verpasst, funktioniert das System nicht.

\textbf{Beispiel – Button-Scan:}

\begin{figure}[htbp]
\centering
\begin{tikzpicture}[
    phase/.style={draw, fill=InfoBg, minimum height=0.6cm, font=\footnotesize},
    reserve/.style={draw, fill=TipBg, minimum height=0.6cm, font=\footnotesize}
]
% Timeline
\draw[thick] (0,0) -- (12,0);
\node[font=\small\bfseries] at (6,0.8) {\SI{10}{\milli\second} Scan-Intervall};

% Phases (not to scale for visibility)
\node[phase, minimum width=1.5cm] at (0.75,0) {};
\node[font=\footnotesize] at (0.75,-0.5) {CD4021};
\node[font=\tiny, color=MutedText] at (0.75,-0.9) {\SI{20}{\micro\second}};

\node[phase, minimum width=1cm] at (2.25,0) {};
\node[font=\footnotesize] at (2.25,-0.5) {Logic};
\node[font=\tiny, color=MutedText] at (2.25,-0.9) {\SI{5}{\micro\second}};

\node[phase, minimum width=1.5cm] at (3.5,0) {};
\node[font=\footnotesize] at (3.5,-0.5) {HC595};
\node[font=\tiny, color=MutedText] at (3.5,-0.9) {\SI{20}{\micro\second}};

\node[reserve, minimum width=7.5cm] at (8.25,0) {};
\node[font=\footnotesize] at (8.25,-0.5) {Reserve};
\node[font=\tiny, color=MutedText] at (8.25,-0.9) {\SI{9955}{\micro\second}};
\end{tikzpicture}
\caption{Timing-Budget eines Scan-Zyklus}
\label{fig:timing-budget}
\end{figure}

\textbf{Budget bei \SI{10}{\milli\second} Intervall:}

\begin{table}[htbp]
\centering
\begin{tabular}{@{}lrr@{}}
\toprule
\textbf{Operation} & \textbf{Dauer} & \textbf{Anteil} \\
\midrule
CD4021 read (2 Bytes) & $\sim$\SI{20}{\micro\second}   & 0,2\,\% \\
Debounce + Selection  & $\sim$\SI{5}{\micro\second}    & 0,05\,\% \\
HC595 write (2 Bytes) & $\sim$\SI{20}{\micro\second}   & 0,2\,\% \\
\textbf{Summe}        & $\sim$\SI{45}{\micro\second}   & $\sim$0,5\,\% \\
\textbf{Reserve}      & $\sim$\SI{9955}{\micro\second} & 99,5\,\% \\
\bottomrule
\end{tabular}
\caption{Timing-Budget}
\label{tab:timing-grundlagen}
\end{table}

\textbf{Was zerstört Timing?}

\begin{table}[htbp]
\centering
\begin{tabular}{@{}ll@{}}
\toprule
\textbf{Gut} & \textbf{Schlecht} \\
\midrule
\code{vTaskDelayUntil()} (FreeRTOS) & \code{delay()} im Hot-Path \\
Logging bei Zustandsänderung        & Logging pro Zyklus \\
Kurze SPI-Transfers                 & Blockierende UART-Ausgabe \\
\bottomrule
\end{tabular}
\caption{Timing-Killer vermeiden}
\label{tab:timing-killer}
\end{table}

\textbf{Rechnung – UART-Overhead:}

Bei \SI{115200}{Baud} dauert ein Zeichen $\frac{10\,\text{Bits}}{115200\,\text{Baud}} \approx \SI{87}{\micro\second}$.
Eine Debug-Zeile mit 50 Zeichen: $50 \times \SI{87}{\micro\second} = \SI{4,35}{\milli\second}$ – fast die Hälfte des Scan-Intervalls!

\begin{warnbox}[Konsequenz]
Debug-Ausgaben nur bei Events (Zustandsänderungen), nie pro Scan-Zyklus.
\end{warnbox}

% -----------------------------------------------------------------------------
\subsection{Debouncing – Mechanisches Prellen filtern}

\textbf{Problem:} Mechanische Taster prellen. Ein einzelner Druck erzeugt mehrere Flanken in wenigen Millisekunden.

\begin{lstlisting}[style=wrap,numbers=none]
Ideal:     ----+          +----
               |          |
               +----------+
                 gedrueckt

Real:      ----+ ++ ++    +++----
               | || ||    |||
               +-++-++----++-
                   Prellen
\end{lstlisting}

\textbf{Messwerte:}

\begin{table}[htbp]
\centering
\begin{tabular}{@{}lr@{}}
\toprule
\textbf{Parameter} & \textbf{Typischer Wert} \\
\midrule
Prellzeit     & 5–\SI{50}{\milli\second} \\
Prellvorgänge & 3–20 Flanken \\
\bottomrule
\end{tabular}
\caption{Typische Prellwerte}
\label{tab:prellwerte}
\end{table}

\textbf{Lösung – Software-Debounce:}

\begin{lstlisting}[style=arduino]
typedef struct {
    uint8_t stable;     // Letzter stabiler Zustand
    uint8_t counter;    // Zaehler fuer aufeinanderfolgende gleiche Samples
} DebounceState;

bool debounce_update(DebounceState* state, uint8_t current, uint8_t threshold) {
    if (current == state->stable) {
        state->counter = 0;  // Zustand bestaetigt, Reset
        return false;
    }

    state->counter++;
    if (state->counter >= threshold) {
        // Neuer Zustand ist stabil
        bool wasLow = (state->stable == 0);
        bool isLow  = (current == 0);
        state->stable = current;
        state->counter = 0;

        // Press-Edge: Active-Low, war HIGH, ist jetzt LOW
        return (wasLow == false && isLow == true);
    }
    return false;  // Noch nicht stabil
}
\end{lstlisting}

\textbf{Berechnung des Threshold:}

Bei \SI{10}{\milli\second} Scan-Intervall und gewünschter \SI{50}{\milli\second} Debounce-Zeit:
\[
\text{threshold} = \frac{\SI{50}{\milli\second}}{\SI{10}{\milli\second}} = 5
\]

\begin{tipbox}[Konsequenz]
Die Logik sieht nur stabile Zustände und saubere Flanken, kein Prellen.
\end{tipbox}

% -----------------------------------------------------------------------------
\subsection{One-Hot-Encoding – Die Invariante}

\textbf{Problem:} Bei einem Selection Panel soll immer genau eine Auswahl aktiv sein (oder keine). Wie stellst du das sicher?

\textbf{One-Hot-Prinzip:}

\begin{lstlisting}[style=wrap,numbers=none]
selectedId = 0:  00000000 00000000  (nichts ausgewaehlt)
selectedId = 1:  10000000 00000000  (ID 1 aktiv)
selectedId = 2:  01000000 00000000  (ID 2 aktiv)
selectedId = 9:  00000000 10000000  (ID 9 aktiv)
\end{lstlisting}

\textbf{Invariante:} Zu jedem Zeitpunkt ist höchstens ein Bit gesetzt.

\textbf{Implementierung:}

\begin{lstlisting}[style=arduino]
void setOneHot(uint8_t selectedId, uint8_t* ledState, size_t bytes) {
    // 1. ALLE Bits loeschen (garantiert sauberen Zustand)
    memset(ledState, 0x00, bytes);

    // 2. Genau EIN Bit setzen (wenn ID > 0)
    if (selectedId >= 1 && selectedId <= LED_COUNT) {
        uint8_t byteIdx = (selectedId - 1) / 8;
        uint8_t bitPos  = 7 - ((selectedId - 1) % 8);  // MSB-first
        ledState[byteIdx] |= (1u << bitPos);
    }

    // 3. Ungueltige Bits im letzten Byte maskieren
    ledState[bytes - 1] &= LED_LAST_MASK;
}
\end{lstlisting}

\begin{infobox}[Warum \code{memset} zuerst?]
Andere Ansätze wie \glqq altes Bit löschen, neues setzen\grqq{} können bei Bugs oder Race Conditions zu Zuständen mit zwei gesetzten Bits führen. Mit \code{memset} + einem \code{|=} ist die Invariante garantiert.
\end{infobox}

% -----------------------------------------------------------------------------
\subsection{SSOT – Single Source of Truth}

\textbf{Problem:} Wenn Pin-Definitionen in drei Dateien stehen und du eine änderst, sind die anderen falsch.

\textbf{Regel:} Jede Information hat genau eine autoritative Quelle.

\textbf{Im Projekt:}

\begin{table}[htbp]
\centering
\begin{tabular}{@{}ll@{}}
\toprule
\textbf{Aspekt} & \textbf{SSOT-Datei} \\
\midrule
Pin-Zuordnung & \code{firmware/include/hw/board\_pins.h} \\
Kanal-Anzahl  & \code{firmware/include/hw/io\_counts.h} \\
SPI-Parameter & \code{firmware/include/hw/spi\_params.h} \\
\bottomrule
\end{tabular}
\caption{SSOT-Dateien im Projekt}
\label{tab:ssot-grundlagen}
\end{table}

\textbf{Beispiel –} \code{io\_counts.h}:

\begin{lstlisting}[style=arduino]
#define BTN_COUNT  10
#define LED_COUNT  10

// Abgeleitete Werte (automatisch korrekt bei Aenderung von COUNT)
#define BTN_BYTES  ((BTN_COUNT + 7) / 8)
#define LED_BYTES  ((LED_COUNT + 7) / 8)

// Maske fuer ungueltige Bits im letzten Byte
#define LED_LAST_MASK  0xC0  // Bei 10 LEDs: nur Bit 7, 6 gueltig
\end{lstlisting}

\begin{tipbox}[Konsequenz]
Änderung von 10 auf 100 Kanäle erfordert nur zwei Zeilen in einer Datei.
\end{tipbox}

% -----------------------------------------------------------------------------
\subsection{Test-Strategie – Vertrauen durch Verifikation}

\textbf{Problem:} Embedded-Bugs sind schwer zu debuggen. Du siehst nicht, was auf der Leitung passiert.

\textbf{Lösung:} Zweistufige Test-Strategie.

\paragraph{Unit-Tests (Logik isoliert)}

\begin{lstlisting}[style=shell]
pio test -d firmware -e seeed_xiao_esp32s3_test -v
\end{lstlisting}

\textbf{Was wird getestet:}
\begin{itemize}
\item Debounce-Algorithmus (filtert Prellen korrekt?)
\item Selection-Logik (One-Hot-Invariante garantiert?)
\item Bit-Mapping (ID → Byte/Bit korrekt?)
\end{itemize}

\textbf{Vorteil:} Läuft in Sekunden, keine Hardware nötig.

\paragraph{E2E-Tests (Hardware verifiziert)}

\begin{lstlisting}[style=shell]
pio test -d firmware -e seeed_xiao_esp32s3_e2e -f test_e2e_hw -v
\end{lstlisting}

\textbf{Was wird getestet:}
\begin{itemize}
\item SPI-Timing (Daten kommen korrekt an?)
\item Pin-Zuordnung (Button 1 → Bit 7 von Byte 0?)
\item LED-Mapping (selectedId=5 → LED 5 leuchtet?)
\end{itemize}

\textbf{Ablauf (interaktiv):}
\begin{enumerate}
\item Idle-Test: Nichts drücken, RAW sollte 0xFF sein
\item Button 1 drücken: RAW Bit 7 = 0?
\item Button 10 drücken: RAW Byte 1, Bit 6 = 0?
\item Walking-One: LEDs durchlaufen, visuell prüfen
\end{enumerate}

\begin{warnbox}[Konsequenz]
Nach Änderungen an SSOT-Dateien oder Verkabelung müssen beide Test-Stufen grün sein.
\end{warnbox}

% -----------------------------------------------------------------------------
\subsection{Zusammenfassung – Merkregeln}

\paragraph{Spannungspegel}

\begin{table}[htbp]
\centering
\begin{tabular}{@{}ll@{}}
\toprule
\textbf{Zustand} & \textbf{\SI{3,3}{\volt} Logik} \\
\midrule
LOW  & \SI{0}{\volt} – \SI{0,8}{\volt} \\
HIGH & \SI{2,0}{\volt} – \SI{3,3}{\volt} \\
\bottomrule
\end{tabular}
\caption{Spannungspegel-Referenz}
\label{tab:pegel-referenz}
\end{table}

\paragraph{Active-Low/High}

\begin{table}[htbp]
\centering
\begin{tabular}{@{}lll@{}}
\toprule
\textbf{Komponente} & \textbf{Typ} & \textbf{Aktiv =} \\
\midrule
Buttons (mit Pullup) & Active-Low  & 0 \\
LEDs                 & Active-High & 1 \\
\bottomrule
\end{tabular}
\caption{Active-Low/High im Projekt}
\label{tab:active-referenz}
\end{table}

\paragraph{SPI-Modi (im Projekt)}

\begin{table}[htbp]
\centering
\begin{tabular}{@{}lcll@{}}
\toprule
\textbf{IC} & \textbf{Mode} & \textbf{CLK Idle} & \textbf{Sample bei} \\
\midrule
CD4021  & 1 & LOW & Falling Edge \\
74HC595 & 0 & LOW & Rising Edge \\
\bottomrule
\end{tabular}
\caption{SPI-Modi-Referenz}
\label{tab:spi-referenz}
\end{table}

\paragraph{Timing-Budget}

\begin{table}[htbp]
\centering
\begin{tabular}{@{}rrr@{}}
\toprule
\textbf{Scan-Intervall} & \textbf{IO-Operationen} & \textbf{Reserve} \\
\midrule
\SI{10000}{\micro\second} & $\sim$\SI{45}{\micro\second} & $\sim$\SI{9955}{\micro\second} \\
\bottomrule
\end{tabular}
\caption{Timing-Budget-Referenz}
\label{tab:timing-referenz}
\end{table}

% -----------------------------------------------------------------------------
\subsection{Weiterführende Dokumentation}

\begin{description}
\item[\code{firmware/docs/wiring.md}] Hardware-Verdrahtung
\item[\code{firmware/docs/architecture.md}] Software-Architektur
\item[\code{docs/grundlagen-bit-operationen.md}] Bit-Manipulation
\item[\code{docs/bit-mapping-schieberegister.md}] ID-zu-Bit-Zuordnung
\end{description}
